% Source: http://tex.stackexchange.com/a/5374/23931
\documentclass{article}
\usepackage[T1]{fontenc}
\usepackage{parskip}
\usepackage{microtype}
\usepackage[utf8]{inputenc}
\usepackage[margin=1in]{geometry}
\pagestyle{plain}
\usepackage{hyperref} 
\setlength{\parindent}{0.5cm}
\newcommand{\HRule}{\rule{\linewidth}{0.5mm}}

\makeatletter% since there's an at-sign (@) in the command name
\renewcommand{\@maketitle}{%
  \parindent=0pt% don't indent paragraphs in the title block
  \centering
  {\LARGE \bfseries\textsc{\@title}}
  \HRule\par%
  \textbf{Gawtam Chithra Ramesh \hfill Computer Science}
  \par
  \textbf{PhD Applicant \hfill AI Core, CARIAD, Munich}
  \par
}
\makeatother% resets the meaning of the at-sign (@)

\title{Motivation Letter}

\begin{document}
  \maketitle
\vspace{7pt}


I am applying to the AI-Core PhD position at CARIAD because end-to-end learning for driving is exactly the direction I have been drawn to for years: a single model that learns perception, prediction, and planning jointly, and improves in closed loop against a world model rather than only imitating a log. Reinforcement learning has been the thread running through my work since 2020, and this role is one of the clearest places I have seen to take it to one of the hardest real-world problems in the field.

That thread started at IIT Madras. As part of the Computer Vision and Intelligence Group, I co-organized a deep learning summer school for over 500 students covering deep learning, RL, and computer vision, and co-mentored student projects on reinforcement learning for games. Around the same time I built a multi-agent RL shooter in Unity3D, training cooperative PPO policies that learned to outplay human opponents, and integrated it into VR. That project is where I fell for the idea that an agent can learn behaviour from its own experience instead of being told what to do, and I have been chasing that idea ever since.

I kept pulling on it through my master's at KTH. For a Capture-the-Flag environment I built a hybrid agent that paired a custom Behavior Tree for high-level strategy with role-specific PPO policies for attack, defend, and escape, trained with self-play and curriculum learning through Unity ML-Agents in a partially observable, multi-agent setting. That project taught me both the power and the brittleness of RL in the open world: reward design, distribution shift, and the gap between imitation and genuine closed-loop learning, which is exactly the gap an end-to-end driving stack has to close. It also taught me how I like to work, which I confirmed at Ericsson Research, where we set out only to build and understand the problem, and two papers came out of five months anyway. I want to fall in love with the problem and let the publications follow.

My more recent work has pushed me toward the generative and world-model side that this role is built on. My master's thesis at ABB Robotics develops a generative flow-matching trajectory planner, a learned motion prior you sample and steer at inference time, which has convinced me that generative models and world models are how end-to-end planning gets past the ceiling of pure imitation. Alongside it I reimplemented one-step flow matching, masked autoencoders, and contrastive learning from scratch in PyTorch and JAX, and deployed and stress-tested a vision-language-action policy on unseen tasks. I have also built the vehicle-grounded planning and control to match: Conflict-Based Search for 20+ holonomic and non-holonomic agents, hybrid A* with kinematic constraints, RRT* for Dubins cars, and heading and sliding-mode controllers for an autonomous ground vehicle from its dynamics.

What I would add beyond the technical fit is how I work. Every project I have described was independent research: my supervisors set the direction and gave feedback, but the brainstorming, the implementation, and the evaluation were mine to own. The way I make progress is to understand an idea by building it, the quickest prototype that shows whether it actually holds, because re-implementing a method teaches me things that reading about it never does. And I have learned to commit to concrete but ambitious goals rather than broad ones: a sharp target is a better springboard into a wider set of questions than a method chosen because it merely feels general. Climbing the wrong hill only tires you, but changing hills constantly means never climbing at all. That is the mindset I want to bring to prototyping, training, and evaluating end-to-end driving models on real datasets and benchmarks at CARIAD, and to publishing and discussing the results with an interdisciplinary team building the next generation of automated driving.

\end{document}